\documentclass[11pt,a4paper]{article}
\newcommand{\doctag}{Requirements review}
\usepackage[a4paper,margin=2.25cm,headheight=15pt]{geometry}
\usepackage[T1]{fontenc}
\usepackage{lmodern,microtype,graphicx,booktabs,array,tabularx,amsmath}
\usepackage[table]{xcolor}
\usepackage{enumitem,listings,fancyhdr}
\usepackage[hidelinks]{hyperref}
\definecolor{navy}{HTML}{163A5F}
\definecolor{pale}{HTML}{EDF3F8}
\definecolor{muted}{HTML}{526272}
\setlength{\parindent}{0pt}
\setlength{\parskip}{0.48em}
\setlength{\emergencystretch}{2em}
\renewcommand{\arraystretch}{1.16}
\newcolumntype{Y}{>{\raggedright\arraybackslash}X}
\newcolumntype{P}[1]{>{\raggedright\arraybackslash}p{#1}}
\setlist{nosep,leftmargin=1.5em,topsep=0.35em}
\pagestyle{fancy}
\fancyhf{}
\fancyhead[L]{\small ERTS Assignment 2\quad\textcolor{muted}{\doctag}}
\fancyhead[R]{\small 11 October 2026}
\fancyfoot[R]{\small\thepage}
\renewcommand{\headrulewidth}{0.4pt}
\lstset{basicstyle=\ttfamily\footnotesize,breaklines=true,columns=fullflexible,
  keepspaces=true,showstringspaces=false,frame=single,rulecolor=\color{muted!30},
  backgroundcolor=\color{pale},xleftmargin=0.2em,xrightmargin=0.2em}
\newcommand{\banner}[1]{\colorbox{pale}{\parbox{\dimexpr\linewidth-2\fboxsep\relax}{#1}}}
\newcommand{\docopening}[2]{
 {\Large\bfseries\color{navy} Embedded Real Time Systems\par}
 \vspace{0.35em}{\LARGE\bfseries #1\par}
 \vspace{0.4em}{\large #2\par}
 \vspace{0.5em}Firas Harbo Saleh \quad Student ID: 202109174\par
 \textcolor{muted}{Review snapshot: \texttt{assignment\_2/Firas}, commit \texttt{c38a216}.}\par
 \vspace{0.6em}}
\newcommand{\src}[1]{{\small\path{#1}}}

\begin{document}
\docopening{Requirements Review}{Mandatory Exercises 3, 4, 5 and 7}

\section*{Post-review closure: 11 October 2026}
This summary covers local fixes on \texttt{assignment\_2/Firas} after the
reviewed commit. The original PDFs and LaTeX sources were preserved under
\src{outputs/fix_evidence/pre_fix_documents/}.

\begin{tabularx}{\linewidth}{P{3.0cm} Y}
\toprule Review finding & Implemented correction and evidence \\
\midrule
Ex3 LED/GPIO mismatch & Custom LED register-0 writes and semantic SDT/legacy
switch lookup. Native build and physical UART pass. \\
Ex3 OCM requirement & Canonical linker script explicitly selected; each actual
ARM ELF has heap \texttt{0x1000--0x2FFF}, stack \texttt{0x3000--0x67FF};
driver configuration section retained. \\
Ex3 counter & UART records every value 0--15. First-run arrival spacing is
0.992--1.008 seconds. User confirmed the full physical LED count during rerun. \\
Ex4 interfaces/fixture & Course \texttt{vectorType}/\texttt{VectorArray} and
function signatures restored; literal transpose gives repeated row values
30/70/110/150. Host tests and native Ex4 run pass. \\
Ex5 missing measurements & Three real-board comparisons pass both equality and
known matrix values. Software ticks: 264, 197, 192; hardware: 3190, 3116, 3122. \\
Ex7 reference discrepancy & Original kept unchanged; explicitly adapted course
reference added with four additional signed/zero/extreme cases. Literal
untouched-testbench acceptance is still unresolved. \\
Benchmark/build robustness & Application \texttt{-O2}; shared ANN dimensions,
driver/address guards, padded transfers, full transaction timing and ten trials.
Workspaces preserved by rename and board cable selected explicitly. \\
\bottomrule
\end{tabularx}

The Ex345 native bitstream met timing (WNS $+1.216$ ns, WHS $+0.046$ ns,
zero negative slack) and final DRC reported zero violations. Original DDR
skew warnings remain visible in the build log rather than being suppressed.

\subsection*{Remaining acceptance boundaries}
The course specifies 2017.2; native success in 2026.1 does not establish course
approval. The original ANN test encodes overlapping row strides and differs
from the implemented contiguous fixed-point model. Approval of the documented
adapted reference is required if literal provided-testbench compliance is graded.
The confirmed LED counter does not by itself prove every switch-to-LED mapping.

\newpage
\section*{Fresh verification and evidence}
The expanded ANN flow has passed both native C simulations, baseline and
pipelined synthesis, and all five VHDL co-simulation transactions. HLS latency
remains 33,056 / 6,612 cycles. Co-simulation latency is 6,610 / 6,612 / 6,613
cycles (min/average/max); the adapted course vector is included. Whole-system
Vivado implementation was rebuilt after export and again met timing with
WNS $+3.462$ ns, WHS $+0.023$ ns and zero negative slack.

\subsection*{Corrected optimized benchmark}
The fresh Vitis build verifies \texttt{-O2 -g3} in the actual compiler command.
The initial and ten subsequent ZYBO transactions all returned matching custom
vector outputs and overall PASS. Repeated measurements use the same inputs:

\begin{center}
\begin{tabular}{lrrr}
\toprule Measurement & Min ticks & Mean ticks & Max ticks \\
\midrule
CPU reference & 13,844 & 14,373 & 14,757 \\
AXI loading & 228,396 & 228,451 & 228,567 \\
HW start-to-done & 21,623 & 21,632 & 21,636 \\
HW FULL transaction & 253,315 & 253,386 & 253,559 \\
\bottomrule
\end{tabular}
\end{center}
At 325 MHz the means are 44.2 / 702.9 / 66.6 / 779.6 $\mu$s. The optimized
ARM is $1.51\times$ faster than hardware computation alone; FULL hardware is
$17.63\times$ slower than CPU compute. The earlier $8.61\times$ comparison
against \texttt{-O0} is superseded, not retained as the current conclusion.

FULL includes packing, writes, readiness, start/done, readback and unpacking;
it excludes platform/data setup, verification and printing. Timer overhead and
cache effects remain; ten repeated samples are not a broad workload study.

\subsection*{Evidence locations}
\src{outputs/fix_evidence/board_ex3_20261011/},
\src{board_ex3_counter_repeat_20261011/}, \src{board_ex4_20261011/} and
\src{board_ex5_20261011/} contain real UART/XSDB logs and artifact hashes.
The rerun folder also retains the user's physical counter observation.
Native compiler/OCM evidence is in
\src{Assignment_2/hardware/Ex345_vitis/native_build_evidence.json} and each app's
\texttt{build/ocm\_verification.json}, ELF, map and compiler-command database.

New ANN evidence: \src{outputs/fix_evidence/ex7_hls_fresh.log},
\src{ex7_vivado.log}, \src{ex7_vitis_retry.log}, and
\src{board_ex7_20261011/} with UART/XSDB logs and bitstream/ELF hashes.
Compiler evidence is also at
\src{Assignment_2/hardware/Ex7_vitis/neuralNetTest/build/build_configuration.txt}.

\newpage
\section*{Historical review: pre-fix snapshot}
The following sections preserve the findings at commit \texttt{c38a216} before
the authorized fixes. Their words ``current'', ``missing'' and ``outstanding''
refer to that snapshot. The closure summary on the preceding pages gives the
updated state and supersedes those status claims without erasing the audit trail.

\banner{\textbf{Verdict: not yet ready for an unqualified hand-in.}
Exercise 7 has genuine native build and physical-board evidence, but several
written requirements remain unmet or unverified. The earlier statement that
all code-level deliverables were complete was too broad.}

\section*{Scope and method}
This review checks the original assignment (25 August 2025, pp.~2--12),
the installation guide, supplied project/IP archives, current sources,
generated reports, and actual ZYBO UART output. Exercises 1, 2 and 6 provide
supporting labs; the mandatory journal covers 3, 4, 5 and 7. Audio playback
is a separate side project, not a missing mandatory deliverable.

Findings distinguish defects, missing evidence and ambiguous course notation.
Application code was not changed. The 11 October board retry reused the existing
native bitstream and ELF; it was not a clean rebuild or another Vivado session.

\section*{Priority findings}
\begin{enumerate}
\item \textbf{Exercise 3 targets the wrong LED peripheral.} The supplied
platform uses custom \texttt{led\_ip\_0}; the application assumes an AXI GPIO
called \texttt{leds}. This prevents using the application unchanged on the
intended hardware, even in the older SDK flow.
\item \textbf{Exercise 3 memory placement and modern GPIO support are incomplete.}
No linker/build artifact proves heap and stack in on-chip RAM (OCM).
Its GPIO initialization still uses legacy device IDs under a modern SDT flow.
\item \textbf{Exercises 4--5 are not demonstrated on the target.} The software
fixture passes, but the prescribed matrix interfaces differ and the printed
transpose notation needs resolution. Exercise 5's required measured
software/hardware comparison is absent.
\item \textbf{Exercise 7 does not pass the supplied testbench.} The custom
testbench passes; the untouched course testbench fails against the same core.
The discrepancy stems from different weight-row indexing and needs an explicit,
accepted explanation, not an undocumented replacement.
\item \textbf{Acceptance and reporting remain open.} Physical switch/LED
observation is missing. The course explicitly specifies 2017.2, while this
work uses 2026.1. Performance claims need compiler and timing-boundary qualifiers.
\end{enumerate}

\section*{What is already supported}
The matrix fixture and custom ANN host/HLS/RTL tests pass. Native implementation
met timing. Physical ZYBO runs on 10 and 11 October produced ten matching
software/hardware outputs. These results do not prove the other requirements.

\newpage
\section*{1. Exercises 3 and 4: requirements and gaps}
\small
\begin{tabularx}{\linewidth}{P{2.8cm} P{2.6cm} Y}
\toprule Requirement & Finding & Evidence / closure condition \\
\midrule
Ex3 UART interpreter, commands 1 and 2 & Source logic exists & Parser, switch read and timed 0--15 counter exist; target integration is blocked by LED-driver mismatch. \\
Ex3 correct LED device & \textbf{Mismatch} & \texttt{main.c:21--28} uses \texttt{XGpio}. Supplied BD has custom \texttt{led\_ip\_0} at \texttt{0x43C00000}. \\
Ex3 one-second increments & Unverified physically & Uses 325,000,000 timer ticks. Observe the LED sequence and actual timing on the board. \\
Ex3 heap/stack in OCM & \textbf{Missing proof} & Supply the linker script, built ELF/map and section addresses, not just a stated intention. \\
Ex4 global matrices & Present & \texttt{aInst}, \texttt{bTInst} and \texttt{pInst} exist. \\
Ex4 prescribed types and functions & \textbf{Different interfaces} & Course image specifies \texttt{vectorType} and \texttt{VectorArray}; current type names and initializer/display signatures differ. \\
Ex4 values and row-dot-row multiply & Partial / ambiguous & Current fixture passes modulo-256 arithmetic; its input orientation differs from a literal reading of p.~6. \\
Ex4 \texttt{mult}, UART result and ticks & Source only & No retained Ex4 native build or physical UART measurement. \\
\bottomrule
\end{tabularx}
\normalsize

\subsection*{The LED and GPIO issues are separate}
The supplied Ex3 archive's \texttt{system\_bd.tcl} instantiates
\texttt{xilinx.com:user:led\_ip:1.0}. Course lab 4 uses a custom LED-IP register
write, and its RTL drives the LEDs from register 0. Replacing the device-ID
macro alone will not turn that device into AXI GPIO.

For switches, the installed 2026.1 \texttt{xgpio.h} selects a base-address
initialization argument under \texttt{SDT}. Current Ex3 code passes
\texttt{XPAR\_SWITCHES\_DEVICE\_ID}; the shared timer already has a conditional
legacy/SDT path, but GPIO does not. A compatible fix must respect the actual
platform's generated address map.

\subsection*{Matrix orientation: resolve before changing expectations}
Let $C$ denote the array stored in \texttt{bTInst}. The implementation computes
$P_{ij}=\sum_k A_{ik}C_{jk}\pmod{256}$, i.e.\ $P=AC^T$.
The current $C$ stores constant rows $1,2,3,4$, yielding rows
$(10,20,30,40)$, $(26,52,78,104)$, $(42,84,126,168)$ and
$(58,116,174,232)$.

The course p.~6 prints a transpose outside that constant-row matrix.
Literal transposition gives four stored rows $(1,2,3,4)$ and results with
repeated row values $30,70,110,150$. Page 7's notation is also loose.
This is an interpretation issue, not evidence that the row-dot-row arithmetic
is wrong. The hand-in must state its chosen stored layout and use consistent
matrix labels; its current table does not do so.

\newpage
\section*{2. Exercise 5: acceleration is not yet measured}
The supplied IP takes two packed 32-bit rows in registers 0 and 1 and returns
an eight-bit dot product in register 2. Truncating each product before summing
and truncating the final sum is equivalent to the software result modulo 256.
The current \texttt{multiMatrixHard} follows this register protocol.

The actual supplied IP uses \texttt{S00\_AXI}; the course example uses
\texttt{S\_AXI}. The code's \texttt{S00\_AXI} spelling matches the supplied
package and is not a defect. However, source-level agreement is not a native
platform build, a physical result, or a measured comparison.

\textbf{Required closure:} build the matrix-IP hardware and application;
capture the full software and hardware matrices, equality result and both
tick counts over UART. Convert ticks using 325 MHz and discuss the 650 MHz
CPU versus 100 MHz PL clocks and the cost of AXI accesses. The existing
journal's measurement plan cannot substitute for these results.

\section*{3. Exercise 7: all nine numbered items}
\small
\begin{tabularx}{\linewidth}{P{0.7cm} P{3.9cm} Y}
\toprule Item & Requirement & Current evidence / verdict \\
\midrule
7.1 & C++ IP and testbench & Implemented; custom self-checking reference passes. \\
7.2 & Simulation/synthesis with provided testbench & \textbf{Unmet literally.} Native reports use the replacement testbench. Untouched supplied testbench fails a direct host check. \\
7.3 & AXI4-Lite arrays & Implemented. Generated control and array buses are both connected and addressed. \\
7.4 & Pipeline both inner loops; compare latency and FF/LUT/DSP/BRAM & Supported by baseline/pipelined reports. Both use 2026.1 and a 10 ns target. \\
7.5 & Lowest-latency C/RTL co-simulation & VHDL PASS, 6,610 cycles for the custom test. One transaction does not measure transaction initiation interval. \\
7.6 & Export VHDL, version 1.0.1 & Successful export command requests 1.0.1. XML stores version 1.0 plus revision 1; this alone is not a failure. \\
7.7 & Vivado IP integration and pin constraints & Saved BD, XDC, native implementation and physical numeric run support completion for this ZYBO. \\
7.8 & SDK application verification and physical switch/LED test & Numeric check passes in Vitis-built application. SDK substitution is a version deviation. Physical LED observation remains unverified. \\
7.9 & Software/hardware timing; separate loading and compute; bottleneck & Measured, with limits: ARM build is \texttt{-O0}; load-plus-run subtotal excludes packing and output readback. \\
\bottomrule
\end{tabularx}
\normalsize

\textbf{Additional model requirements:} fixed-point arithmetic, Q15 shifts,
hidden-layer ReLU and switch-to-LED assignment are present. Dimensions are
compile-time macros; the default 196/32/10 and a temporary 8/3/4 host variant
pass the custom test. Only the default configuration has synthesis/RTL/board
evidence. The SDK application is independently fixed at the default dimensions.

\newpage
\section*{4. Neural-network testing: reconcile the references}
The untouched provided \texttt{Ex7ANN/testbench.cpp} and current accelerator
were compiled together with the installed AMD arbitrary-precision headers.
That executable exits 1: all ten comparisons fail. Two representative values:

\begin{center}
\begin{tabular}{lrr}
\toprule Output & Current accelerator & Supplied floating reference \\
\midrule
0 & 331 & 41.003418 \\
9 & 4805 & 201.729141 \\
\bottomrule
\end{tabular}
\end{center}

The original reference starts input-weight rows at \texttt{NEURONS\_LEN*i}
and output-weight rows at \texttt{OUTPUT\_LEN*i}. The implemented fully
connected model instead uses \texttt{INPUT\_LEN*i} and
\texttt{NEURONS\_LEN*i}. Those are sensible contiguous row strides; reverting
them blindly would overlap rows. The supplied reference also uses floating
arithmetic, whereas the implemented reference matches the prescribed integer
shifts. These are modeling/test changes, not consequences of a newer IDE.

\textbf{Recommended resolution, not an implemented fix:} retain the supplied
test as provenance, document the row-stride discrepancy, and add a clearly
labeled adapted course reference with justified fixed-point expectations.
Run that test through C simulation and RTL co-simulation, and seek course
acceptance for the deviation. Passing only the new synthetic fixture does not
prove the literal supplied-testbench requirement.

\section*{5. Verified native results}
\begin{center}
\begin{tabular}{lrrrrr}
\toprule HLS solution & Cycles & BRAM18K & DSP & FF & LUT \\
\midrule
Baseline & 33,056 & 10 & 2 & 638 & 716 \\
Pipelined inner loops & 6,612 & 10 & 3 & 1,112 & 1,126 \\
\bottomrule
\end{tabular}
\end{center}
The approximately $5\times$ latency reduction costs one DSP, 474 FF and
410 LUT, with unchanged BRAM18K usage. This compares two solutions in the
\emph{same} 2026.1 toolchain, not Vivado 2017.2 versus 2026.1.

VHDL co-simulation reports 6,610 cycles. Routed system timing reports WNS
$+3.462$ ns, TNS 0, WHS $+0.023$ ns and THS 0. Routed system usage is
959 LUT, 1,097 registers, six BRAM tiles and three DSPs. These are whole-system
implementation figures, not directly interchangeable with HLS core estimates
or the BRAM18K unit count.

The fresh 11 October 2026 physical run downloaded the existing bitstream and
ELF through the existing hardware server and returned:
\[
[-4,\ 3,\ -2,\ 1,\ -4,\ 2,\ -1,\ 0,\ -3,\ 0].
\]
Every hardware value matched the software reference; UART ended with
\texttt{Overall verification: PASS}. The generated VHDL also contains a direct
\texttt{outLeds <= inSwitch} connection. This is continuous pass-through,
but does not itself prove that someone observed the physical LEDs.

\newpage
\section*{6. Timing boundaries and reporting corrections}
\begin{center}
\begin{tabular}{lrrr}
\toprule Measurement & 10 Oct ticks & 11 Oct ticks & 11 Oct time \\
\midrule
ARM software compute & 186,751 & 186,747 & 574.6 $\mu$s \\
AXI input/weight writes & 233,860 & 233,849 & 719.5 $\mu$s \\
Accelerator start-to-done & 21,697 & 21,700 & 66.8 $\mu$s \\
Measured load-plus-run subtotal & 255,557 & 255,549 & 786.3 $\mu$s \\
\bottomrule
\end{tabular}
\end{center}
Time is ticks divided by 325 MHz. The fresh compute-only ratio is
$186747/21700\approx8.61$. The application compile command uses
\texttt{-O0 -g3}; this is a comparison with that unoptimized ARM reference,
not a best-possible CPU implementation. Two recorded runs are useful
repeat observations, not a statistically characterized benchmark.

Input packing occurs before the loading timer. Hardware execution timing starts
before the start-register write and includes done polling; output reads and
unpacking happen after it stops. The subtotal therefore excludes preparation,
readback and setup. It must not be called full end-to-end inference latency.
Even this partial subtotal exceeds software compute, so input/weight traffic
is already a demonstrated bottleneck for the current one-shot workflow.

\subsection*{Corrections needed in the existing hand-in}
\begin{itemize}
\item Replace ``code complete'' / ``ready for supplied Ex3 platform'' with the
actual peripheral and linker gaps until corrected and built.
\item Limit Ex7's all-checks-passed statement to the custom tests; disclose the
provided-testbench failure and required physical LED test.
\item Correct matrix labels and explicitly document the stored orientation.
\item Qualify the $8.61\times$ number as compute-only versus an
\texttt{-O0} reference; relabel transfer-plus-compute as a subtotal.
\item State the 2017.2 requirement deviation. The course does not merely
recommend that version, and acceptance of 2026.1 has not been established.
\end{itemize}
The original hand-in and code were left unchanged by this review. This document
is a review addendum, not a claim that those corrections have already been made.

\subsection*{Minimum acceptance checklist before submission}
\begin{enumerate}
\item Resolve tool-version and supplied-reference deviations with the course.
\item Fix Ex3 against the actual custom LED IP, use compatible GPIO lookup,
and prove OCM heap/stack placement from the built application.
\item Align the matrix interfaces and clarify inputs; capture Ex3--5 native
UART results, physical LED/counter observations and measured Ex5 comparison.
\item Run the documented course-derived neural test and capture physical
switch/LED observations. Suggested coverage: patterns 0, 1, 2, 4, 8 and 15.
\item Correct the hand-in's result claims and retain logs with build identity.
\end{enumerate}
These are outstanding assignment actions. Completion of this review is not
completion of the assignment itself.

\newpage
\section*{7. Evidence index and reproduction}
Paths below are relative to the workspace
\src{/home/firas/Documents/Codex/2026-10-08/i-h/}.
Repository references use prefix \src{ERTS_assigments/Assignment_2/}.

\small
\begin{tabularx}{\linewidth}{P{2.65cm} Y}
\toprule Evidence & Location / what it establishes \\
\midrule
Course specification & \src{work/assignment2/Assignment_2/ERTS_Assignment2_ANN_ZYBO.pdf}; pp.~4--9 and 10--12 are the exercise requirements. \\
Installation baseline & \src{/home/firas/Downloads/ERTS_Vivado_Installation_2017_2.pdf}; pp.~1 and 4 specify 2017.2 and SDK. \\
Ex3 source / platform & \src{exercise3_console/src/main.c}; original \src{/home/firas/Downloads/Ex3_ZYBO.xpr.zip}, including \texttt{system\_bd.tcl} and custom LED RTL. \\
Matrix implementation & \src{common/matrix.h}, \src{common/matrix.c}, \src{exercise4_matrix/src/main.c}, \src{exercise5_matrix_ip/src/main.c}, \src{tests/test_matrix.c}. \\
Original ANN reference & \src{work/assignment2/starter/Assignment_2/Ex7ANN/testbench.cpp}; supplied archive extraction, not the replacement in Git. \\
HLS implementation & \src{exercise7_neural_net/hls/}; source, header, replacement testbench and \texttt{run\_hls.tcl}. \\
Native HLS reports & Below that HLS directory: \src{neuralNet_hls/baseline/syn/report/neuralNet_csynth.rpt}, corresponding \texttt{pipelined} report, and \src{neuralNet_hls/pipelined/sim/report/neuralNet_cosim.rpt}. \\
Vivado integration & \src{hardware/Ex7/Ex7.srcs/sources_1/bd/system/system.bd}; constraints under \src{hardware/Ex7/Ex7.srcs/constrs_1/imports/Ex7ANN/contraints.xdc}. \\
Timing / area & \src{hardware/ex7_timing_summary.rpt}, \src{hardware/ex7_utilization.rpt}. \\
Application boundaries & \src{exercise7_neural_net/sdk/neuralNetTest.c:104--153}; packing/loading/run/readback; flags in \src{hardware/Ex7_vitis/neuralNetTest/build/compile_commands.json}. \\
Physical retry & \src{outputs/native_screenshots/zybo_retry_20261011T085514Z/uart.log} and \texttt{xsdb.log}; fresh hardware evidence. Blank screenshot attempt is excluded. \\
Original recorded run & \src{outputs/native_screenshots/ex7_board_uart.log}; 10 October benchmark. \\
\bottomrule
\end{tabularx}
\normalsize

The matrix check was rerun as \texttt{make -C Assignment\_2 test} from the
repository root and passed. Lightweight ANN checks use \texttt{g++ -std=c++14}
with \texttt{-I/home/firas/AMD/2026.1/Vitis/include}. Compile the current core
once with the replacement testbench and separately with the untouched course
testbench; they exit 0 and 1 respectively. The temporary 8/3/4 check modifies
only a copied header, not repository files.

Retained host output is in \src{outputs/review_evidence/}: the matrix fixture,
default custom ANN, alternate 8/3/4 ANN and untouched-course ANN logs.
\texttt{alternate\_neuralNet.h} preserves the temporary dimension settings.

Generated bitstreams, BSPs, ELF files and HLS outputs are largely ignored by
Git. Their local presence is evidence for this review, not proof that a clean
clone reproduces them. The companion migration journal records the rebuild
order, provenance and limitations.
\end{document}
